%% problem_id: S048-PJM-305-1-p01-thm-1
%% source_id: S048 (MSP; Pacific Journal of Mathematics)
%% source_article: Kenneth L. Baker, "The Poincare homology sphere, lens space surgeries, and some
%%   knots with tunnel number two", Pacific J. Math. 305 (2020), no. 1, 1--30.
%% source_url: https://msp.org/pjm/2020/305-1/pjm-v305-n1-p01-s.pdf
%% representation: PDF; statement and proof transcribed by RENDERING the pages (pdftoppm -r 150 -png)
%%   and READING THE IMAGES, then checking every display against them.  (The text layer of this paper
%%   additionally misleads: its "Theorem 1.1/1.2/1.4" strings are CITATIONS to other authors, not this
%%   paper's results, whose own numbering is single-digit.)
%% statement_locator: PDF p. 2 (printed p. 1), Theorem 1
%% proof_locator: PDF pp. 8--10 (printed pp. 7--9), "Proof of Theorem 1" through the closing square
%% difficulty_level: H3
%% why H3: the proof is SELF-CONTAINED in this paper and carries a genuine three-dimensional topology
%%   argument -- pass to double branched covers, identify the two tangle pieces and their Seifert
%%   fibrations with exactly two exceptional fibers, exclude the degenerate values of n by exhibiting
%%   compressible tori / tunnel number one, handle the Klein-bottle (twisted I-bundle) case, and close
%%   by the slope-of-a-regular-fiber criterion together with the two fillings that produce two-bridge
%%   pretzel links.  It cites Kobayashi's classification and the Montesinos trick as EXTERNAL standard
%%   results, which plan 2.2 permits; the substance is not delegated (contrast: the NUMDAM paper, whose
%%   Theorem 2.1 is a one-paragraph reduction to another proposition).
%% status: complete (statement + author proof verbatim + reason + locators)
%% ---------------------------------------------------------------------------
\documentclass[11pt]{article}
\usepackage{amsmath,amssymb,amsthm}
\newtheorem{theorem}{Theorem}
\begin{document}

\section*{Theorem 1}

\textbf{Source.} K.~L.~Baker, \emph{The Poincar\'e homology sphere, lens space surgeries, and some
knots with tunnel number two}, Pacific J.\ Math.\ \textbf{305} (2020), no.~1; printed pages 1 and
7--9.

\begin{theorem}
For each integer $n$, the knot $K_n$ in the Poincar\'e homology sphere $\mathcal{P}$ has an integral
surgery to the lens space $L(3n^2+n+1,\,-3n+2)$. If $|n|\ge 4$, then $K_n$ has tunnel number $2$.
\end{theorem}

\subsection*{Proof (authors' text, printed pp.~7--9)}

Figure 1 exhibits arcs $\kappa_n$ on the pretzel knot $P(-2,3,5)$ that have a banding to a two-bridge
link. By passing to the double branched covers (and using the Montesinos trick), this describes knots
$K_n$ in $\mathcal{P}$ that have an integral surgery to a lens space.

Since $K_n$ lies in a genus 2 Heegaard surface, its tunnel number is at most $2$. If $K_n$ were to
have tunnel number one, then every surgery on $K_n$ would have Heegaard genus at most $2$. However we
will see that surgery along the framing induced by the Heegaard surface produces a toroidal manifold
which, according to Kobayashi's classification [1984], does not have Heegaard genus 2.

Kobayashi shows that if $M=M^\alpha\cup_T M^\beta$ is a genus 2 manifold decomposed along a torus $T$
into two atoroidal manifolds, then one of $M^\alpha$ or $M^\beta$ admits a Seifert fibering over the
disk with 2 or 3 exceptional fibers or over the M\"obius band with up to 2 exceptional fibers such
that filling the other along the slope induced by a regular fiber in $T$ produces a lens space. (This
is more general than Kobayashi's classification but suitable for our needs.) The slope of a regular
fiber may be identified since filling one of these Seifert fibered spaces produces a reducible
manifold if and only if the filling is done along the slope of a regular fiber.

Figure 2 (left) shows that the result of a 0-framed banding along $\kappa_n$ has a sphere $S$ dividing
it into two 2-string tangles. Up to homeomorphism, these two tangles are $T_n^\alpha = 1/(n+2)+1/3$
and $T_n^\beta = 1/(1-n)-1/2$. Their corresponding double branched covers, $M_n^\alpha$ and
$M_n^\beta$, are in general each Seifert fiber spaces over the disk with exactly two exceptional (and
nondegenerate) fibers: $M_n^\alpha$ has type $D^2(|n+2|,3)$ and $M_n^\beta$ has type $D^2(|1-n|,2)$.
This fails for $M_n^\alpha$ when $n=-1,-2,-3$ and for $M_n^\beta$ when $n=0,1,2$; in each, the middle
value yields a degenerate Seifert fibration while the other values yield solid tori. In these cases
the torus $T$ that is the double branched cover of $S$ is compressible and so Kobayashi's
classification does not apply; moreover, one may observe that $K_n$ has tunnel number one in these
cases. Hence we assume $n\notin\{-3,-2,-1,0,1,2\}$.

When $n=-1$ or $3$, $M_n^\beta$ has type $D^2(2,2)$ and thus is a twisted $I$-bundle over the Klein
bottle. Therefore it has an alternative Seifert fibration over the M\"obius band with no exceptional
fibers. We already assume $n\neq -1$. We shall assume $n\neq 3$ as well.

Figure 2 (middle) shows the results of filling $T_n^\alpha$ and $T_n^\beta$ with the rational tangle
$\rho^\alpha$ defined by a pair of arcs in $S$ so that $T_n^\alpha(\rho^\alpha)$ is composite. We then
observe that the filling $T_n^\beta(\rho^\alpha)$ is the pretzel link $P(1-n,2,2)$, and this is a
two-bridge link if and only if $n=0,2$. We have already omitted these values of $n$.

Figure 2 (right) shows the results of filling $T_n^\alpha$ and $T_n^\beta$ with the rational tangle
$\rho^\beta$ defined by a pair of arcs in $S$ so that $T_n^\beta(\rho^\beta)$ is composite. We then
observe that the filling $T_n^\alpha(\rho^\beta)$ is the pretzel link $P(n+2,3,-2)$, and this is a
two-bridge link if and only if $n=-1,-3$. We have already omitted these values of $n$ too.

Therefore, to conclude that the knot $K_n$ has tunnel number 2, it is sufficient to require that
$n\notin\{-3,-2,-1,0,1,2,3\}$. \hfill$\square$

\end{document}
